%chapter 2
\chapter{Analysis of RLC parallel circuit changing the R, L and C
values}

We analyzed the RLC parallel circuit showed in figure \ref{fig:ltspice1} by changing the values of R, L and C as following: 
\begin{table}[h]
\centering
\begin{tabular}{c @{\hspace*{1cm}} c c @{\hspace*{1cm}} c c}
\hline
 & L (\si{\nano\henry}) & R (\si{\kilo\ohm}) & C (\si{\femto\farad}) \\
\hline
\text{[A]} & 30 & 200 & 50 \\
\text{[B]} & 30 & 50 & 50 \\
\text{[C]} & 60 & 100 & 50 \\
\text{[D]} & 15 & 100 & 50 \\
\text{[E]} & 30 & 100 & 100 \\
\text{[F]} & 30 & 100 & 25 \\
\hline
\end{tabular}
\caption{Values of L, R and C for different cases}
\end{table}


We measured the resonance frequency $f_r$ and quality factor $Q$ for each case using LTSpice and compared them with the analytical results.
\begin{table}[h]
\centering
\begin{tabular}{c @{\hspace*{1cm}} c c @{\hspace*{1cm}} c c}
\hline
 & $f_r$ analytical (GHz) & $f_r$ experimental (GHz) & $Q$ analytical & $Q$ experimental \\
\hline
\text{[A]} & 4.109 & 4.109 & 258.198 & 257.810 \\
\text{[B]} & 4.109 & 4.109 & 64.549  & 64.521  \\
\text{[C]} & 2.905 & 2.906 & 91.287  & 91.233  \\
\text{[D]} & 5.811 & 5.811 & 182.574 & 182.432 \\
\text{[E]} & 2.905 & 2.906 & 182.574 & 182.243 \\
\text{[F]} & 5.811 & 5.811 & 91.287  & 91.243  \\
\hline
\end{tabular}
\caption{Comparison between analytical and experimental resonance frequency and quality factor}
\end{table}

The analytical predictions are consistent with the experimental results.